\documentclass{article}

\usepackage[margin=3cm]{geometry}
\usepackage{polyglossia}
\usepackage{microtype}
\usepackage{amsmath}

\setmainlanguage[numerals=maghrib]{arabic}
\setotherlanguage{english}

\newfontfamily\arabicfont[Script=Arabic,Scale=1.2]{Estedad}
\newfontfamily\englishfont{JetBrainsMono Nerd Font}

\usepackage{bidi}

\newcommand{\en}[1]{\textenglish{#1}}

\title{\en{Travelling Salesman Problem (TSP)}}
\author{\en{Ahmed Al-khlil \& Anas Al-abdolkarem \& Osama mohamed}}
\date{\en{\today}}

\emergencystretch=1.5em
\spaceskip=1.3\fontdimen2\font plus 1.1\fontdimen3\font minus 1.1\fontdimen4\font

\begin{document}
\maketitle

\section{السيناريو}
تخيل أن هناك شركة توصيل لوجستيه تمتلك مركبه توصيل واحدة, وانطلاقاً من المقر الرئيسي للشركة,
يتعين على السائق زيارة عدد محدد من المدن (او العملاء) لتسليم البضائع, ثم العوده في نهايه المطاف
الى نقطه البدايه (المقر الرئيسي).

ولتحقيق اعلى مستوى من الكفاءة وتقليل التكاليف الاضافيه, يواجه السائق عده شروط صارمه:
\begin{enumerate}
	\item   يجب زياره كل مدينه مره واحده فقط.
	\item يجب تقليل اجمالي المسافه المقطوعه (او اجمالي زمن الرحله او الوقود) الى الحد الادنى.
	\item يجب العوده الى نقطه الانطلاق (المقر الرئيسي).
	\item يمنع بتاتا اغلاق حلقه جزئيه ببعض المدن ثم العوده الى نقطه البدايه والخروج مره اخرى لاكمال باقي المدن.
\end{enumerate}

\subsection{التعقيد الحسابي للمشكلة}
تكمن الصعوبه البالغه في هذه المشكله عندما يبدأ عدد المدن بالارتفاع, حيث يزداد عدد المسارات الممكنه
بشكل "عاملي" ($n!$). لنستعرض كيف تتضخم الاحتمالات:

\begin{itemize}
	\item اذا كان على السائق زياره 3 مدن فقط, فان حساب المسار الاقصر سهل للغايه لان الاحتمالات محدوده. 
	\item اذا ارتفع العدد الى 10 مدن, يقفز عدد المسارات المحتمله الى اكثر من 180000 مسار بديل. 
	\item  اما اذا كان هناك 20 مدينه فقط, فان عدد الاحتمالات ينفجر ليصل الى رقم خيالي من\\ الكوينتيليونات ($10^{18}$ مسار). 
\end{itemize}

بسبب هذا النمو الأسي الهائل في الاحتمالات, يصبح من المستحيل على الحواسيب العاديه حساب واختبار كل مسار على حده للوصول الى الحل المثالي مع زياده جحم البيانات. لذلك, تصنف "مشكلة البائع المتجول"  (\en{TSP})  كواحده من أشهر المسائل غير الحتميه متعدده الحدود الصعبه جداً (\en{NP-hard}), وهي المسائل التي يسهل يصاغتها وفهمها, ولكن يصعب جداً حلها بشكل مطلق عندما تصبح واسعه النطاق.

\subsection{انواع المشكله وعدد الوصلات}
تنقسم الى نوعين رئيسيين بتاءاً على اتجاه الطرق:

\begin{enumerate}
	\item \textbf {مشكله البائع الجوال المتناظره \en{Symmetric TSP}}: هنا تكون الطرق ذات اتجاهين ومتساويه في التكلفه. اي ان المسافه من المدينه  A الى المدينه B, تساوي المسافه من A الى B. اي
	($c_{ij} = c_{ji}$). ويحسب عدد الوصلات بالعلاقه:
	{\LARGE $$
	Number of Edges = \frac{n(n - 1)}{2} \quad ; \quad n \in N^*
	$$}
	
	\item \textbf{ مشكله البائع الجوال غير المتناظره \en{Asymmetric TSP}}: وهنا تكون الطرق موجهة وتختلف تكلفاها حسب الاتجاه ($c_{ij} \not = c_{ji}$). ويحسب عدد الوصلات بالعلاقه: 
	{\LARGE $$
	Number of Edges = {n(n - 1)} \quad ; \quad n \in N^*
	$$}
\end{enumerate}

\subsection{طرق حساب وتمثيل البيانات}
عند حساب المسافات بين المدن وتجهيزها قبل ادخالها الى الخوارزميه كي لا تستهلك قدره المعالج في تكرار حسابات غير مهمه, ولحساب المسافات بين اي نقطتين هناك عده قوانين رياضيه, منها قانون حساب المسافه الاقليديه وهي المسافه المباشره بين نقطتين, وتعطى بالعلاقه التاليه:

{\LARGE
$$
d_{ij} = \sqrt{(x_i - x_j)^2 + (y_i - y_j)^2}
$$
}
\vspace{2mm}

ولتمثيل البيانات رياضيا او برمجياً  يتم وضع البيانات في مصفوفه ثتائيه الابعاد  2D  حجمها {\large$n \times n$}.

\begin{itemize}
	\item السطر يعبر عن مدينه الانطلاق, والعمود يعبر عن مدينه الوصول.
	\item التقاطع بين السطر $i$ والعمود يعبر عن قيمه التكلفه $c_{ij}$ (المسافه او الزمن).
	\item القيم على القطر الرئيسي تكون دائما صفراً او مالانهايه, لمنع السائق من الانتقال من المدينه الى المدينه نفسها.
\end{itemize}

\subsection{اختيار الخوارزميه المناسبه}
\subsubsection{معايير المقارنه في اختيار الخوارزميه المناسبه}
عند التعامل مع مشلكه البائع المتجول \en{TSP}. لا يوجد مفهوم "الخوارزميه الافضل مطلقاً" دون النظر الى سياق البيانات وحجم المشكله. 
ان عمليه اختيار الخوارزميه المناسبه تحكمها دائماً مقايضه اساسيه (\en{Tread-off}) 
بين ثلاثة عوامل رئيسيه وهي: الزمن المستغرق (\en{Time}),
الموارد المتاحه (اي القدره الحسابيه والتخزينيه) (\en{Calculation \& Memory}), 
والدقه المطلويه (\en{Accuracy}). \\

بناءً على هذه الفلسفه, ولتحديد الخوارزميه التي تحقق التوازن الامثل بين كفائه الاداء ودقه النتائج,
قمنا باجراء مقارنه ممنهجه بين سته من ابرز الخوارزميات عالمياً, تتراوح بين الحلول الجشعه البسيطه, والبرمجه الدبناميكيه المعقده, وصولاً الى الخوارزميات الهجينه المتطوره. وفقاً لخمسه معايير محوريه:

\begin{enumerate}
	\item فئه الخوارزميه (\en{Algorithm class}): التصنيف البنيوي للخوارزميه.
	\item فلسفه اتخاذ القرار  (\en{Descision logic}): الاليه التى تتبعها الخوارزميه لاستكشاف المسارات.
	\item التعقيد الحسابي والزمني (\en{Time \& Space complexity}): مدى استهلاك الخوارزميه للوقت وموادر النظام مع نمو عدد المدن ($n$).
	\item دقه الحل (\en{Accuracy}): مدى قرب مسار الناتج من الحل الاقصر المثالي.
	\item مرونه البيانات (\en{Data flexibilty}): القدره على التعامل مع القيود المختلفه والمسارات موجهة الاتجاه.
\end{enumerate}

\subsubsection{مقارنه الخوارزميات}

\begin{itemize}
	\item \textbf{الجار لاقرب (\en{Nearest Neighbor})}:
		\begin{enumerate}
			\item فئه الخوارزميه (\en{Algorithm class}): جشعه (\en{Gready Heuristic}).
			\item فلسفه اتخاذ القرار  (\en{Descision logic}): اختيار المدينه الاقرب محلياً في كل خطوه دون النظر للمستقبل.
			\item دقه الحل (\en{Accuracy}): سيئه الى متوسطه (غالباً اطول بحوالي 25\% من الحل المثالي).
			\item التعقيد الحسابي والزمني (\en{Time \& Space complexity}): التعقيد الزمني اسي $O(n)$, والذاكره $O(n^2)$ او $O(n)$.
		\end{enumerate}

	\item \textbf{هليد-كارب (\en{Held-Karp})}:
		\begin{enumerate}
			\item فئه الخوارزميه (\en{Algorithm class}): دقه فائقه (\en{Exact/Dynamic Programming}).
			\item فلسفه اتخاذ القرار  (\en{Descision logic}): تقسيم المشكله الى مشاكل فرعيه ابسط وتخزين حلولها لمنع تكرار العمليات الحسابيه وحساب الحل المثالي المطلق.
			\item دقه الحل (\en{Accuracy}): الدقه مثاليه 100\%, ولكن لاتستطيع معالجه اعداد مدن تزيد عن 32 مدينه تقريباً.
			\item التعقيد الحسابي والزمني (\en{Time \& Space complexity}): التعقيد الزمني انفجاري 
				$O(n^22^n)$, والتعقيد التخزيني $O(2^nn)$ هائل جداً.
		\end{enumerate}
	\item \textbf{مستعمره النمل (\en{Ant Colony Optimization})}:
		\begin{enumerate}
			\item فئه الخوارزميه (\en{Algorithm class}): ذكاء السراب (\en{Metaheuristic}).
			\item فلسفه اتخاذ القرار  (\en{Descision logic}): محاكاه سلوك النمل عبر افراز الفيرومون, والمسارات الاقصر تكتسب جذباً اقوى تدريجياً عبر ترسب الفيرومون.
			\item دقه الحل (\en{Accuracy}): عاليه جداً (تقترب من المثاليه في الاحجام المتوسطه).
			\item التعقيد الحسابي والزمني (\en{Time \& Space complexity}): التعقيد الزمني $O(I \times A \times n^2)$, حيث I التكرارات وA عدد النمل. 
			والتعقيد التخزيني $O(n^2)$.
		\end{enumerate}
	\item \textbf{التطور الجيني (\en{Generational Optimization})}:
		\begin{enumerate}
			\item فئه الخوارزميه (\en{Algorithm class}): جينيه تطوريه (\en{Metaheuristic}).
			\item فلسفه اتخاذ القرار  (\en{Descision logic}): توليد مجتمع من الحلول, ثم اجراء تزاوج (\en{Crossover}), وطفرات (\en{Mutation}) للبقاء للاصلح والاقصر.
			\item دقه الحل (\en{Accuracy}): منوسطه الى عاليه (تحتاج وقت طويل للتقارب).
			\item التعقيد الحسابي والزمني (\en{Time \& Space complexity}): التعقيد الزمني $O(I \times P \times n^2)$, حيث I التكرارات وP حجم المجتمع. 
			والتعقيد التخزيني $O(nP)$.
		\end{enumerate}
	\item \textbf{\en{opt-2 \& opt-3}}:
		\begin{enumerate}
			\item فئه الخوارزميه (\en{Algorithm class}): تحسين محلي (\en{Local Search}).
			\item فلسفه اتخاذ القرار  (\en{Descision logic}): قطع وتوصيل الوصلات بين المدن,
			حافتين في كل مره (\en{2-opt}),
			او ثلاثه في كل مره (\en{3-opt}). وذلك لايجاد وصلات اقصر
			\item دقه الحل (\en{Accuracy}): متوسطه الى جيده جداً
			\item التعقيد الحسابي والزمني (\en{Time \& Space complexity}): 
			العقيد الزمني $O(n^2)$ \en{for opt-2}, 
			و $O(n^3)$ \en{for opt-3}.
		\end{enumerate}
	\item \textbf{LKH-3}:
		\begin{enumerate}
			\item فئه الخوارزميه (\en{Algorithm class}): تحسين محلي متطور (\en{Advanced Metaheuristic}).
			\item فلسفه اتخاذ القرار  (\en{Descision logic}): 
			تقطيع وتوصيل للوصلات بشكل ديناميكي متغير تابع ل K من الوصلات (opt-k),
			مع ميزات ذكيه لمنع السقوط في الحلول المحليه الفرعيه وتدعم القيود المعقده.
			\item دقه الحل (\en{Accuracy}): قريبه جداً من المثاليه (99\% الى 100\%).
			\item التعقيد الحسابي والزمني (\en{Time \& Space complexity}):
			في التعقيد الزمني, تقترب جداً من $O(n^{2.2})$.
			وفي التعقيد التخزيني $O(n)$ ذو كفائه عاليه جدا في اداره الذاكره.
		\end{enumerate}
\end{itemize}

\section{خوارزميه \en{LKH-3}}
\subsection{مدخل الى عبقريه خوارزميه LKH-3}
في عام 1973, قدم العالمان \en{Shen Lin} و \en{Brian Kernighan} خوارزميه LK , والتي ابتكرت فكرة استخدام k-opt متغير. فبدلاً من تبديل عدد ثابت من الوصلات (مثل وصلتين او ثلاث في كل مره), تغير الخوارزميه عدد تبديل الوصلات بناءً على المعطيات ديناميكياً للخروج من فخ الانماط دون المستوى الامثل (\en{Suboptimal Patterns}).\\

في عام 2000, احدث العالم \en{Keld Helsgaun} ثوره في هذا البحث عندما قدم خوارزميه LKH. وقد اضاف محرك رياضي قوي لتوجيه محرك البحث المحلي: \en{The Held-Karp Lower Bound}. وذلك بدمج \en{Subgradient Optimaization} مع محرك تبديل الوصلات \en{LK swapping engine}, بذلك اصبحت خوارزميه \en{LKH} اقوى خوارزميه حل لمشكله \en{TSP} على الاطلاق.\\

في عام 2017, اصدر العالم \en{Helsgaun} خوارزميه \en{LKH-3}, حيث طور عده اصدارات من المشكله, مثل اضافه خاصيه توجيه عده بائعين, نوافذ الوقت, والتكديس, مما جعلها تتراس على عرش الخوارزميات الاكفئ في حل المشكله عند ازدياد حجمها.\\

اسطاعة خوازرميه \en{LKH-3} تحقيق سرعتها بتقسيم المشكله الى اجزاء اصغر مركزه وذات كفائه عاليه.

\subsection{الخطوه الاولى: تهيئه البيانات}
في المرحله الاولى نعد بيانات المدن ونقوم بحساب الوصلات بينها ونخزنها في \en{Arrays} تسهل النعامل معها.

\subsection{الخطوه الثانيه: \en{The Held-Karp optimaization loop}}
\subsubsection{القاعده الرياضيه: \en{Lagrangian Relaxation}}
قبل تبديل اي وصله, تقوم \en{LKH} بتشغيل \en{subgradient optimaization} لحساب معاملات العقوبه او التحول للمسافه (\en{penalized distances}) لكل مدينه. وثم تنشئ بنيه بيانيه تدعى \en{1-Tree} في كل دوره. اذا كانت مدينه معينه مزاره الكثير من المرات (\en{degree > 2}), قيمه ال\en{penalty} خاصتها ستزداد, مما يجعل التوصيل معها مكلف اكثر, اما اذا كانت معزوله (\en{degree < 2}), تنقص قيمه ال\en{penalty} خاصتها ,بعد مئات من هذه الدورات, تؤدي هذه العمليه الى تسطح الرسم البياني, مما يرفع الحد الادنى حتى يصل الى ما دون الحل المثالي للمشكله مباشرةً.\\

يمكن تمثيل مشكله البائع المتجول برقم وحيد: وهو تقليل اجمالي تكلفه الرحله مع مراعاه القيد الذي ينص على ان كل مدينه يحب ان يكون لها وصلتين فقط (\en{degree = 2}) بالضبط دون انشاء حلقات فرعيه منفصله. ولكن بسبب اجبار الدرحات ان تكون 2 بالضبط لكل المدن في وقت واحد هو ما يجعل المشكله من تصنيف (\en{NP-Hard}). استخدم طريقه \en{Held-Karp} تكنيكاً يدعى (\en{Lagrangian Relaxation}). وبدلاً من اجبار قيد الدرجه بالقوه, نقوم بازاله القيود الثابته وننقلها الى داله تحسين الهدف (\en{penalty}). لاي مصوفه من الارقام الحقيقه $\pi = (\pi_1, \pi_2, \dots, \pi_n)$, نعرف مصوفه من قيم $\pi$ المعدله $C'$:
\large$$
	C'_{ij} = C_{ij} + \pi_i + \pi_j
$$
\newpage
حيث $C_{ij}$ هي المسافه الجفرافيه الحيقيقه بين المدينه $i$ والمدينه $j$. والقيم $\pi_i$ و $\pi_j$ هي مضاعفات لاغرانج. اذا جمعنا هذه القيم $C'_{ij}$ لكل مدينه, تصبح التكلفه المعدله لمجموعه مدن تحتوي $N$ وصلات متغيره او تابعه بشكل منهجي تباعاً للدرجات الخاصه بكل مدينه $d_{ij}$:\\

$$
\begin{aligned}
\text{\en{Penalized Cost}} = W(\pi) &= \sum_{e \in E'} C'_{i,j} = \sum_{(i,j) \in E'}  C_{i,j} + \sum_{i=1}^{N} d_i\pi_i \\
									&= \sum_{e \in E'} C'_{i,j} = \sum_{(i,j) \in E'} (C_{i,j} + \pi_i + \pi_j)
\end{aligned}
$$\\

حيث $E'$ هي مجموعه جزئيه من مجموعه الوصلات $E$.

\subsubsection{اليه عمل الحلقه: \en{Subgradient Ascent}}
تبدا عمليه التحسين باستخدام  \en{Subgradient} لايجاد اقصى قيمه للحد الادنى, تعمل الحلقه على اربع خطوات متكرره:

\begin{enumerate}
	\item حساب \en{Minimum 1-Tree}:
	باستخدام قيم الاوزان المعدله (اوزان العقوبه) $C'_{ij}$ الحاليه, يتم تشغيل خوارزميه \en{Kurskal} عليها لايجاد شجره \en{MST} على $N - 1$ مدن, ثم نقوم بوصل المدينه المميزه (المعزوله) معها باستخدام اقصر وصلتين عندها (من حيث قيمه \en{pi} ايضاً). بذلك يكون اجمالي الوزن المعدل الناتج معطى ب $W(\pi)$.\\
	\begin{itemize}
	\item السبب وراء اختيار البنيه \en{1-Tree}:
	يعد إيجاد مسار صحيح لمسألة البائع المتجول \en{TSP} أمرًا بالغ الصعوبة بسبب قيدٍ صارم: يجب أن يكون حلقةً واحدةً متصلة حيث تتصل كل مدينة بمدينتين أخريين فقط. ويُعتبر إيجاد بنيةٍ ذات وزنٍ أدنى وبهذه الخاصية تحديدًا كابوسًا من فئة NP-Hard.

تتجاوز شجرة \en{1-Tree} هذه الصعوبة بتخفيف القواعد بما يكفي لتسهيل العمليات الحسابية، مما ينشئ بنيةً تعد بديلًا مثاليًا لمسار \en{TSP}.

رياضيًا، شجرة \en{1-Tree} هي بنية بيانية مبنية على $N$ مدينة تحتوي على $N$ وصله بالضبط وحلقة مغلقة واحدة فقط، ويتم ذلك باستخدام
حيلة تفكيك بسيطة من خطوتين:\\
		\begin{enumerate}
		\item أنشئ الشجرة الأساسية، واختر مدينة واحدة لتكون "المدينه الخاصة" (مثل المدينة 0). أزلها تمامًا من الخريطة، ثم أنشئ شجرة امتداد دنيا قياسية (\en{MST}) عبر المدن المتبقية N-1 باستخدام خوارزمية \en{Kruskal}.\\

		\item بمجرد العثور على الشجرة الممتدة الدنيا، نعود إلى المدينة 0 ونقوم بتوصيلها قسرًا بالشجرة الممتدة الدنيا الناتجة باستخدام أرخص حافتين متاحتين لها. بناءً على تكاليفك الحالية.\\
		\end{enumerate}

	لأن الشجرة الممتدة الدنيا القياسية لا تحتوي على حلقات، وعقدة إضافية واحدة (المدينة 0) ستضمن إنشاء دورة واحدة بالضبط.\\

	\item لماذا تعد \en{1-Tree} مرآه مثاليه لمشكله \en{TSP}؟\\
		لنلقِ نظرة على المتطلبات الهيكلية لمسار صالح في مسألة البائع المتجول \en{TSP} مقارنة بـ \en{1-Tree}:\\
		\begin{itemize}
		\item تحتوي جولة \en{TSP} الصحيحة على $N$ من الحواف، متصلة بالكامل، ولها دورة واحدة فقط، ولكل عقدة درجة تساوي $2$ بالضبط.\\

		\item تحتوي الشجرة أحادية البعد على $N$ حافة متصلة بالكامل، ولها دورة واحدة فقط، لكن درجات عقدها يمكن أن تكون أي قيمة (درجة $1$ أو أكثر)، وهو عدد الحواف المتصلة بها.\\
		\end{itemize}
		\vspace{1em}
		وهذا يعني أن مسار \en{TSP} الصالح هو ببساطة نوع خاص من بنية\\ \en{1-Tree}، حيث تكون درجة كل مدينه مساوية لـ 2، كما أن قواعده وبنيته البسيطة تجعل عملية بنائه سهلة للغاية.\\

	\item عملية استخراج الشجرة الممتدة الدنيا \en{MST}:\\
		لاستخراج الشجرة الممتدة الدنيا \en{Minimum Spaning Tree (MST)}، نستخدم خوارزمية \en{Kruskal}، بما أننا نبني شجرة أحادية، فإننا نتجاهل مدينه واحده، مما يترك لنا عدد N - 1 مدينة، وفي نظرية الرسوم البيانية (\en{Graph Theory})، تتطلب الشجرة الممتدة الدنيا التي تربط مجموعة من الرؤوس $V$ دائمًا $V - 1$ حافة بالضبط لربطها جميعًا دون تكوين حلقة مغلقه، وبالتالي سيكون إجمالي عدد الحواف $(N-1)-1$ ($1$ المدينه المميزه التي عُزلت).\\
		\item الاداه الاساسيه: \en{The Disjoint Set Union (DSU)}:\\
أثناء قيام خوارزمية \en{Kruskal} بفحص الوصلات المرتبة (من الأرخص إلى الأغلى)، تطرح سؤالًا بسيطًا: "إذا أضفت هذه الوصله، فهل ستُنشئ حلقة مغلقة غير قانونية؟"، يجيب هيكل المجموعة المنفصلة (\en{Disjoint Set}) على هذا السؤال فورًا من خلال تتبع المدن المتصلة بالفعل بفروع هيكلية. ولذلك يستخدم مصفوفتين بسيطتين:\\
	\begin{itemize}
		\item \textbf{\en{parents}}: يتتبع الجذر المباشر (الأصل)   لكل مدينة في مجموعتها، إذا كانت المدينة هي أصلها، فهي جذر مجموعتها المستقلة.\\

		\item \textbf{\en{ranks}}: يتتبع عمق فروع الشجرة لمجموعة ما لضمان أنه عند دمج مجموعتين، يتم توصيل المجموعة الأقصر أسفل المجموعة الأطول.
	\end{itemize}\vspace{1em}

	عند بدء تشغيل الخوارزمية، تكون كل مدينة كيانًا منفردًا ومعزولًا تمامًا؛ إذ لا تعرف أي مدينةٍ منها المدن الأخرى. وتشكل كل مدينة "دائرة اجتماعية" مستقلة خاصة بها، تتألف من عنصر واحد فقط. وتتولى بنية بيانات "المجموعات المنفصلة" \en{Disjoint Set} إدارة هذه الدوائر باستخدام عمليتين أساسيتين:\\

		\begin{itemize}

		\item \textbf{\en{Find}}:
		يسأل مدينة: "من هو القائد المطلق (الأصل الجذري) لدائرتك الاجتماعية؟"\\

		\item \textbf{\en{Union}}:
		يجمع بين دائرتين اجتماعيتين منفصلتين تماماً تحت قيادة واحدة موحدة.
		\end{itemize}\vspace{1em}


أثناء تكرار الخوارزمية على الحواف المرتبة، تأخذ معرفي مدينتين وتتحقق من
جذورهما (كل منهما في دائرتها الاجتماعية أو مجموعتها الخاصة). إذا لم تكن جذورهما متساوية، فإنها تأخذ الحافة بينهما بأمان، ثم تدمج المجموعة ذات الترتيب الأصغر
في المجموعة ذات الترتيب الأعلى، بحيث يكون جذر المجموعة ذات الترتيب الأعلى هو جذر المجموعة ذات الترتيب الأصغر. أما إذا كانت الجذور متساوية، فإنها تتخطى كليهما لأن لهما مسارًا مشتركًا أقصر، وبالتالي لا يحتاجان إلى حافة مباشرة بينهما من شأنها أن
تسبب حلقة مغلقة.\\

	\end{itemize}

	\item حساب \en{Subgradient Vector (G)}:
	يمثل التدرج الفرعي الاتجاه الرياضي لاكبر خطأ, لكل مدينه $i$, نحسب كم تنحرف درجتها الحاليه $d_i$ في شجره \en{1-Tree} عن الهدف المراد في اي رحله \en{TSP} فعليه (حيث كل مدينه درجتها تساوي 2 تماماً). يمكننا تعريف الخطأ الهيكلي الفردي لأي عقدة $i$ على النحو التالي:\\

		{\LARGE$$
		G_i = d_i - 2
	$$}\\

لفهم كيف يُمثل التدرج الفرعي الاتجاه الرياضي لأكبر خطأ، من المفيد الابتعاد عن البرمجة والنظر إلى هندسة التحسين. إذا كنت تتسلق جبلاً في ضباب كثيف وتريد الوصول إلى القمة بأسرع وقت ممكن، فإنك تنظر إلى الأرض تحت قدميك وتخطو في الاتجاه الذي يصعد فيه المنحدر بأكبر انحدار ممكن. في علم التفاضل والتكامل، يُسمى اتجاه الصعود الأكبر
بالتدرج (\en{Gradient - grad}). ومع ذلك، فإن رسم دالة الهدف في تحسين Held-Karp ليس جبلاً أملساً ومستديراً تماماً. لأن الشجرة أحادية البعد تتكون من حواف حادة منفصلة، يبدو رسم الدالة وكأنه ماسة متعددة الأوجه أو مثل الهرم. يحتوي على حواف وزوايا حادة حيث يفشل حساب التفاضل والتكامل القياسي لأنه لا يمكنك حساب مشتقة تقليدية عند نقطة حادة. هنا يأتي دور التدرج الفرعي (\en{Subgradient}). حيث إنه تعميم للتدرج للدوال غير الملساء والمتعرجة.\\

في المرحلة الثانية، يتمثل هدفنا الاول هو جوله صالحه (\en{Valid TSP Tour}). في الجولة الصالحة، يكون لكل مدينة هدف بان تكون درجتها تساوي 2 بالضبط. إذا أنتجت الخوارزمية \en{1-Tree} حيث تكون درجة إحدى المدن 4، فإننا نواجه خطأً هيكليًا. اي ان المدينة مكتظة. إذا كانت درجة إحدى المدن تساوي 1، فإننا نواجه خطأً أيضًا؛ المدينة معزولة.\\

اذا كانت المدينه مزدحمه ($d_i = 4$), يكون تدرجها موجباً ($4 - 2 = +2$). اذا كانت المدينه منعزله او طريق مغلق ($d_i = 1$), يكون تدرجها سالباً ($1 - 2 = -1$). واذا كانت المدينه في حاله دوره \en{TSP} مثاليه ($d_i = 2$), ينهار تدرجها نحو الصفر ($2-2=0$).\\

لنتذكر صيغة اجمالي "التكلفة العقابيه" المعممة الخاصة بنا:\\

{\LARGE$$
	W(\pi) = \sum_{(i,j) \in E'} C_{i,j} + \sum_{i=1}^{N} d_i\pi_i \\
$$}\\

لننظر كيف يؤثر تغيير عقوبة مدينة واحدة $\pi_i$ على إجمالي التكلفة المعاقبة $W(\pi)$. إذا اعتبرنا حواف شجرة \en{1-Tree} المختارة ثابتة مؤقتًا، فيمكننا حساب المشتقة الجزئية للتكلفة بالنسبة لعقوبة تلك المدينة:\\

{\LARGE$$
	\frac{\partial W}{\partial \pi_i} = d_i
$$}\\

يُشير هذا إلى أن معدل التغير الحالي (الميل) لدالة التكلفة على طول المحور $\pi_i$ يساوي تماماً الدرجة الحالية للعقدة، $d_i$. ومع ذلك، فإننا نسعى لتعظيم الحد الأدنى، ولكن يجب علينا القيام بذلك دون السماح "للعقوبات" (\en{penalties}) بتشويه المخطط البياني إلى ما لا نهاية. وتقتضي مسألة التحسين المزدوج (\en{dual optimization problem}) تحقيق توازن في النظام، حيث تقوم خطوة التعديل المستهدفة بضبط الميل بالنسبة للقيد المفروض (أي الدرجة المستهدفة البالغة 2)؛ وعند طرح القيد المستهدف، يصبح متجه الميول الاتجاهية كما يلي:\\


$$
	\nabla W(\pi) \approx
	\begin{bmatrix}
		d_1 - 2 \\
		d_2 - 2 \\
		\vdots  \\
		d_N - 2
	\end{bmatrix}
	=
	\begin{bmatrix}
		G_1    \\
		G_2    \\
		\vdots \\
		G_N
	\end{bmatrix}
	=
	\vec{G}
$$\\

هذا المتجه $\vec{G}$ هو التدرج الفرعي الخاص بنا.\\

يخبرنا \en{Subgradient} بالاتجاه الذي يجب سلكه لتعديل اخطائنا البنيويه وماهي النقاط التي علينا تعديلها. على سبيل المثال, اذا كانت المدينه ذات درجه تساوي $5$ ($G_i = +3$). يكون الميل ايجابي واتجاه \en{subgradient vector} يؤشر بشده نحو الاتجاه الموجب. ومع التقدم في الحلقات, يسبب ذلك ازدياد قيمه $\pi_i$ لتلك المدينه بشكل اقوى. وهذا يجعل الوصلات المتصله بتلك المدينه اكبر بالوزن واغلى مع كل حلقه, مما يجبر خوارزمية \en{Kruskal} على اقتطاع الوصلات عاليه التكلفه منها.\\

وعلى العكس تماماً, اذا كانت المدينه ذات درجه منخفصه مثل $1$ ($G_i = -1$). في هذه الحاله يؤشر متجه \en{subgradient} نحو الاتجاه السالب. ومع الوقت, ستنخفض قيمه $\pi_i$ لتلك المدينه. مما يجعل الوصلات المتصله معها ارخص في التكلفه, ويجبر ذلك خوارزمية \en{Kruskal} على اضافه المزيد من الوصلات لتلك المدينه.



	\item حساب حجم الخطوه \en{Step Size (t)}:
	لايمكننا ببساطه تغيير قيم $\pi$ بشكل اعمى؛ فعل ذلك سيؤدي انحراف النظام بشكل كبير وعشوائي عن القيمه الفضلى. يتم حساب حجم التغير عن طريق استخدام متغير \en{setp-size scalar t}, مقاد \en{By Held and Karp's classic convergence formula} :\\

	{\LARGE$$
	t = \lambda \frac{\text{\en{Target Bound}} - W(\pi)}{\sum_{i=1}^{N}(G_i)^2}
	$$}\\

	حيث:
	\begin{itemize}
		\item \en{Lmabda} $\lambda$:
		هو معامل التخميد التكيفي والذي يبدا حول عدد حقيقي ويتلاشى نحو الصفر.
		\item \en{Target Bound}: هو تقدير للقيمه الحقيقه للسقف (\en{Upper Bound}),
		عالبا مايحسب ابتداءً باستخدام خوارزميه بحث سريعه مثل \en{nearest neighbor}, او تعدل بواسطه الحلقه الخارجيه.

		\item المقام $\sum (G_i)^2$:
		مجموع الاخطاء للتربيع, ذلك كان الشكل العام لتوزيع المدن عشوائياً والدرجات بعيده عن 2, تكبر قيمه المقام, مما يجعل حجم الخطوه اصغر واكثر استقراراً.

	\end{itemize}

	\item تحديث قيم $\pi$:
	اخيراً, بتم تعديل مصفوفه قيم العقوبات (او يتم تعديلها لكل مدينه) من اجل الدوره التاليه:\\
	اذا كانت درجه المدينه عاليه, قيمه $\pi$ خاصتها تزداد. في الدوره التاليه, كل الوصلات المتصله معها تصبح "اغلى", يدفع ذلك خوارزميه \en{Kurskal} على تجنبها. وبالعكس, المدن المنفيه او ذات الوصله الواحده يحصلون على قيمه تعديل ($\pi$) سالبه, مما يجعل وصلاتهم "ارخص" واكثر جذباً للخوارزميه في الدوره التاليه.
\end{enumerate}




\end{document}
